\documentclass[10pt,norsk,a4paper]{article}
\usepackage[utf8]{inputenc}
\usepackage[T1]{fontenc}
\usepackage[norsk]{babel}
\usepackage[cm]{fullpage}
\usepackage{color}
\usepackage{parskip,textcomp,amssymb,graphicx}
\usepackage{pdfpages}
\usepackage[stable]{footmisc}
\usepackage{multicol}


\title{Generalforsamling \\
	Høsten 2017\\[3cm]
	\includegraphics[width=0.5\textwidth]{cyblogoa3.pdf}\\[-.5cm]}
\date{15.\ november 2017}
\author{Cybernetisk Selskab}

% Blank header, samt footer med side x av y
\usepackage{fancyhdr}
\pagestyle{fancy}
\renewcommand{\headrulewidth}{0pt}
\fancyhead{}
\cfoot{Side~\thepage\ av~\pageref{lastpage}}



\begin{document}

\maketitle{}
\newpage
\tableofcontents{}


\section{Valg av møteleder}

\section{Valg av referent}

\section{Valg av protokollunderskrivere}

\section{Valg av tellekorps}

\section{Godkjenning av innkalling}

\section{Godkjenning av dagsorden}

\section{Semesterberetninger}
\subsection{Semesterberetning ved leder}

\textit{Kjære interne, medlemmer og studenter,}
\begin{multicols}{2}

	Dette har vært et bra år for Cybernetisk Selskab. Vi utvidet fadderuka i år, noe som ble en braksuksess. Vi arrangerte storfester de tre siste dagene av fadderuka, og vi hadde fortsatt ikke kapasiteten vi trengte. Selv om det er kjipt med kø, er det veldig postitivt at vi klarer å tiltrekke oss så mange studenter.

	De to andre store arrangementene vi har holdt dette semesteret---dagen@ifi og Ifi-galla---gikk begge uten store problemer. Spesielt dagen@ifi virker det som om at vi har funnet en fin formel på hvordan vi skal arrangere dette sammen med dagen-styret.

	Vi har også dette semesteret startet med en ny type arrangement: \textit{Torsdagsklubben / spritkveld}. Dette er et konsept som har slått bra an. Det er ikke fult hus på disse kveldene, men det kommer mer enn nok folk til at dette er et konsept vi burde fortsette med.

	Personlig har jeg jobbet en del med det pågående prosjektet med å anskaffe nytt kassaapparat. Planen er per dags dato at denne vil bli installert første eller andre uken i desember, og at det blir tatt i bruk etter nyttår. Dette vil medføre enklere økonomirutiner, mindre teknisk kompleksitet, og at de interne i baren får noe nytt å leke med.

	Medlemsmassen vår er på ca.\ det samme som det var i fjor, med 714 semestermedlemmer, noe som vi kan si oss fornøyde med når vi tenker på at vi satte rekord i fjor.

	Helt til slutt vil jeg nevne at Cybernetisk Selskab fyller 50 år i 2019. Jeg har sammen med noen pamper og andre interesserte ifi-studenter startet i det små for å finne ut hva vi ønsker å få til. Vi ønsker å danne et ``CYB-50''-utvalg som skal ta seg av organiseringen av jubileumet, og at det skal bestå av både gamle og nye studenter. Jeg håper at dette er formelt på plass like etter nyttår. Vi har også snakket med Fordelingsutvalget, som er positive til å støtte prosjektet så lenge vi skriver en intensjonsmelding på hva vi ønsker å gjennomføre, samt sørger for at dette er et arrangement for hele instituttet.

	Vi sees over nyåret!

	\textbf{Andreas Hansen},\\
	leder \\
	\date{10.\ november 2017}
\end{multicols}

\newpage


\subsection{Semesterberetning ved kjellermogul}

\begin{multicols}{2}
	Da nærmer atter et semester seg sin slutt, og ikke annet kan egentlig sies enn dette har vært et svært godt år for CYB og Escape.

	I fadderuka i år kjørte vi utvidet område ikke bare på fredagen, men også på onsdagen og torsdagen. Selv om dette hadde sine utfordringer, gikk det bra for å være et pilotprosjekt, og vi endte opp med en markant høyere omsetning for hele uka, samt at mange flere funksjonærer stilte opp enn året før. Årets fadderstyre fokuserte også veldig på å opprettholde en god kommunikasjon med oss, og hjalp til med å sørge for at alle faddere hadde arbeidsoppgaver, der hvor å jobbe i bar eller som vertskap gikk innunder dette.

	Vi har også hatt enda flere arrangementer enn før, hvor nesten alle ukene i semesteret har hatt noe som har foregått. Dette er egentlig noe som bare blir bedre med hvert år, som kanskje er selvsagt. Dagen@ifi gikk også bedre enn noen gang, hvor vi i år økte omsetning på ca. 50 000,- i forhold til i fjor. Det er tydelig at vi med tiden har begynt å få disse arrangementene mer og mer under huden, og har nå ganske mye informasjon om hvordan man faktisk arrangerer de større arrangementene på instituttet. Jeg kan også meddele at renhold på en generell basis er veldig fornøyde med oss når det gjelder så og si alle de store festene vi arrangerer, noe som jo er positivt.

	Torsdagsklubben er også et nytt prosjekt der vi prøver å jevnt ha drinkkvelder i løpet av semesteret, både for å trene opp folk i hvordan å jobbe med drinker (til større arrangementer hvor vi trenger flinke folk) og for å gi et tilbud til de som er glad i drinker framfor annen alkoholholdig væske for å gi rom for et enda større publikum. Selv om---som også Andreas påpeker---at oppmøtet ikke nødvendigvis er så stort, så er det noe jeg tror bare kommer til å bli bedre med tiden hvis vi klarer å fortsette med slike arrangementer. I samme slengen kan det også nevnes at vi har hatt såpass mange arrangementer at det har ført til at det er mange utlånsforespørsler som vi har måttet takke nei til, fordi det rett og slett ikke har vært plass på ukeplanen. På mange måter vil jeg si at det bare er en positiv indikasjon på Escape sin rolle som møte- og knutepunkt.

	Web- og driftgruppene ble slått sammen til X-gruppa, som jeg tror viser at det er---og kommer til å bli---et enda større fokus rundt både det å få inn flere interne i gruppa, og jobb med utvikling av internsystemet, wikien og nettsidene våre.

	For å følge opp det jeg nevnte sist semester om kjelleravtalen, har Cosa Nostra generelt jobbet ganske lite med den nye (foreløpige), utover dokumentet vi sendte til UiO med krav om ønskede forandringer i forrige kjelleravtale. Så vidt jeg vet sikter universitetet sentralt fremdeles på å fullføre den nye avtalen innen første halvdel av 2018.

	Til slutt vil jeg egentlig bare si takk for meg, og takk for et supert år med to autonome og produktive kjellerstyrer, samt en fantastisk internbase som er de som faktisk får ting til å gå rundt. Jeg er---som sikkert er sagt av mange før meg---veldig positiv til hva CYB kan få til framover, og har fått til på ganske få år, og håper at vi en dag klarer å bli ikke bare en av de største, men den største kjellerforeningen ved Blindern.

	Vi sees i Escape!

	\textbf{Karl Hole Totland}, \\
	kjellermogul \\
	\date{14.\ november 2017}
\end{multicols}

%\newpage

\section{Kasserer orienterer}
\subsection{Regnskapets tilstand}
Kasserer orientere om regnskapet.


\subsection{Revidert budsjett for høsten 2017}
Kasserer orientere om budsjett.



\section{Kontigentfastsettelse}
Hovedstyret foreslår å endre medlemskontigenten fra kr. 30,- til kr. 40,-,
fra og med neste semester. \\
Hovedstyret orienterer rundt motivasjonen bak dette.


\section{Valg}

\begin{minipage}[t]{9cm}
\subsection{Hovedstyret}
Man velges inn i hovedstyret for ett år av gangen.

\subsubsection{Kjellermogul}
\subsubsection{Kasserer}
\subsubsection{Arrangementssjef}
\subsubsection{Promoteringssjef}
\subsubsection{Internansvarlig}

\end{minipage}
\begin{minipage}[t]{9cm}
\subsection{Kjellerstyret}
Alle verv--\textit{utenom økonomiansvarlig}--som er til valg i \\kjellerstyret gjelder for ett semester av gangen.

\subsubsection{Økonomiansvarlig}
\subsubsection{Innkjøpsansvarlig}
\subsubsection{Barsjef}
\subsubsection{Kafésjef}
\subsubsection{Teknisk ansvarlig}
\subsubsection{DJ-sjef}
\subsubsection{Utlånsansvarlig}
\subsubsection{Arrangementskoordinator\footnotemark}

\end{minipage}
\footnotetext{Tentativ vervtittel.}

\newpage

\section{Vedtektsendringer}
Følgende forslag til vedtektsendringer ligger fremme.

\subsection{Endring i forbindelse med krav om Ifi-foreninger}

For å være en offisiell Ifi-forening har Ifi satt sammen noen vedtekter som studentforeningene må følge.
Per dags dato følger ikke Cybernetisk Selskab §1c\footnotemark. Det ble lagt frem et forslag av hovedstyret på vårens generalforsamling om å endre våre vedtekter, slik at vi ikke bryter med Ifi sine krav.
Dette ble da ikke vedtatt da forsamlingen konkludere med at de nye vedtektene var for åpne og uklare. Forsamlingen ønsket at hovedstyret skulle komme med et nytt forslag til neste generalforsamling, som blir denne, høsten 2017.

\footnotetext{Etter manglende vedtektsfesting av to-tredelers Ifi-studenter i hovedstyret}

Om vedtektene ikke blir vedtatt vil hovedstyret starte søknadsprosessen for å få et unntak fra §1c av \textit{Vedtekter for Ifi-forening}, slik som §3c åpner for. Begge endringsforslagene hører sammen og stemmes over i lag.

Hovedstyret er innstilt på å starte søknadsprosessen etter §3c.

\subsubsection{Forslag til ny paragraf: $§5h$}
\begin{quote}
	\begin{enumerate}
		\item[§5h] Minimum to-tredeler av hovedstyret skal bestå av studenter som enten (1) er tilknyttet et eller flere studieprogrammer ved \textit{Institutt for informatikk}, eller (2) som tar ett eller flere IN- eller INF-fag.
	\end{enumerate}
\end{quote}

\subsubsection{Forslag til endring av paragraf: §8a}
\begin{quote}
	\begin{enumerate}
		\item[§8a]
			Valgbare til alle styrer er alle medlemmer i henhold til §2 og informatikkstudenter ved Universitetet i Oslo. \textcolor{red}{Om valg av kandidat bryter med §5h, blir valget gjort ugyldig og det blir gjenvalg. Om ingen andre kandidater stiller blir vervet etterfylt av hovedstyret i henhold til §5f.} Alle valg avgjøres \textcolor{red}{på generalforsamling} ved simpelt flertall.
	\end{enumerate}
\end{quote}


\section{Æresmedlemskap}

\section{Utdeling av pins}
Arkivar deler ut pins.

\section*{Vedlegg: Cybernetisk Selskabs vedtekter}
\section*{Vedlegg: Vedtekter for Ifi-foreninger}\label{lastpage}


\newpage

\includepdf[pages=-]{cyb_vedtekter.pdf}
\includepdf[pages=-]{ifi_forening_vedtekter_2017.pdf}
\end{document}
